(1) 一个量子比特初态 $\rho=(I+x\sigma_x+y\sigma_y+z\sigma_z)/2$。定义二结果非投影测量的效应 $E_\pm=(I\pm\eta\sigma_z)/2$，$0\le\eta\le1$。证明效应的正性和归一性，求测量结果概率。

(2) 取Lüders更新对应Kraus算符 $K_\pm=\sqrt{E_\pm}$，求给定结果后的 $z$ 分量及横向相干振幅变化。

(3) 若不记录结果，求量子通道 $\mathcal E(\rho)=\sum_\pm K_\pm\rho K_\pm$，证明其保持布居并使横向Bloch矢量缩小因子 $\sqrt{1-\eta^2}$。

(4) 讨论 $\eta=0$、$\eta=1$ 的极限，并量化对两种输入本征态 $|\uparrow\rangle,|\downarrow\rangle$ 的测量区分成功率。
